\documentclass[a4paper]{article}

\usepackage[spanish]{babel}
\usepackage{graphicx}
\usepackage{amsmath, amssymb, mathrsfs}
\usepackage[margin=2cm]{geometry}
\usepackage{fancyhdr}
\usepackage{enumerate}
\usepackage[shortlabels]{enumitem}
\usepackage{parskip}
\usepackage[most]{tcolorbox}
\usepackage[hidelinks]{hyperref}
\usepackage{float}
\usepackage{booktabs}



% cabecera
\pagestyle{fancy}
\fancyhead[l]{Física Matemática I}
\fancyhead[c]{Taller 7}

\renewcommand{\headrulewidth}{0.2pt} % linea horizontal

\begin{document}

\begin{titlepage}
  \centering
  {\Large Universidad de Antioquia\par}
  \vspace{2cm}
  {\huge\bfseries Física Matemática I\par}
  \vspace{0.4cm}
  {\Large Solución al Taller 7 \par}
  \vspace{2.5cm}
  {\large Autor:\par}
  \vspace{0.2cm}
  {\large Daniel Soto Villada\par}
  {\large Estudiante de Astronomía\par}
  \vfill
  {\large \today\par}
\end{titlepage}

\newpage
\tableofcontents
\newpage


\section{Problema 1}
Utilizando la función generatriz, muestre que $\frac{2n}{x}J_n(x) = J_{n+1}(x) + J_{n-1}(x)$ y $2J_n'(x) = J_{n+1}(x) - J_{n-1}(x)$.

\subsection{Solución}
Para demostrar las relaciones de recurrencia solicitadas, partimos de la definición de la función generatriz de las funciones de Bessel de orden entero, herramienta fundamental en la teoría de funciones especiales que se introduce y explota en la Sección 8.3 del libro \textit{Lecciones de física matemática} de Alonso Sepúlveda. La función generatriz se define como:
$$g(x, t) = e^{\frac{x}{2}\left(t - \frac{1}{t}\right)} = \sum_{n=-\infty}^{\infty} J_n(x) t^n$$
donde $J_n(x)$ representa las funciones de Bessel de primera especie de orden $n$, y la serie converge para todo $t \neq 0$. Esta expansión permite relacionar funciones de diferente orden mediante operaciones algebraicas y diferenciales sobre el parámetro auxiliar $t$ y la variable física $x$.

\textbf{1. Derivación de la relación $\frac{2n}{x}J_n(x) = J_{n+1}(x) + J_{n-1}(x)$}

Diferenciamos la función generatriz con respecto a la variable auxiliar $t$:
$$\frac{\partial g}{\partial t} = \frac{\partial}{\partial t} \left[ e^{\frac{x}{2}\left(t - \frac{1}{t}\right)} \right] = \frac{x}{2}\left(1 + \frac{1}{t^2}\right) e^{\frac{x}{2}\left(t - \frac{1}{t}\right)} = \frac{x}{2}\left(1 + \frac{1}{t^2}\right) \sum_{n=-\infty}^{\infty} J_n(x) t^n$$
Por otro lado, diferenciando la expansión en serie término a término respecto a $t$:
$$\frac{\partial g}{\partial t} = \sum_{n=-\infty}^{\infty} n J_n(x) t^{n-1}$$
Igualando ambas expresiones y distribuyendo el factor $\frac{x}{2}$ sobre la serie:
$$\sum_{n=-\infty}^{\infty} n J_n(x) t^{n-1} = \frac{x}{2} \sum_{n=-\infty}^{\infty} J_n(x) t^n + \frac{x}{2} \sum_{n=-\infty}^{\infty} J_n(x) t^{n-2}$$
Para poder comparar coeficientes, es necesario alinear las potencias de $t$ en las tres sumatorias a la forma $t^{n-1}$. Realizamos los siguientes cambios de índice discreto:
- En la primera sumatoria del lado derecho, sustituimos $n \to n-1$, obteniendo: $\sum J_n(x) t^n = \sum J_{n-1}(x) t^{n-1}$.
- En la segunda sumatoria del lado derecho, sustituimos $n \to n+1$, obteniendo: $\sum J_n(x) t^{n-2} = \sum J_{n+1}(x) t^{n-1}$.
Sustituyendo estos ajustes en la igualdad:
$$\sum_{n=-\infty}^{\infty} n J_n(x) t^{n-1} = \frac{x}{2} \sum_{n=-\infty}^{\infty} J_{n-1}(x) t^{n-1} + \frac{x}{2} \sum_{n=-\infty}^{\infty} J_{n+1}(x) t^{n-1}$$
Dado que esta igualdad debe cumplirse para todo $t \neq 0$ en el dominio de convergencia, la independencia lineal de las potencias de $t$ exige que los coeficientes de $t^{n-1}$ sean idénticos término a término:
$$n J_n(x) = \frac{x}{2} J_{n-1}(x) + \frac{x}{2} J_{n+1}(x)$$
Reorganizando algebraicamente y despejando el cociente, obtenemos la primera relación de recurrencia:
$$\frac{2n}{x} J_n(x) = J_{n-1}(x) + J_{n+1}(x)$$
Este resultado coincide exactamente con la identidad de recurrencia listada en la Sección 8.3.1 de Alonso Sepúlveda, y demuestra cómo la simetría $t \leftrightarrow 1/t$ en la función generatriz induce una relación de simetría entre órdenes adyacentes.

\textbf{2. Derivación de la relación $2J_n'(x) = J_{n-1}(x) - J_{n+1}(x)$}

Ahora diferenciamos la función generatriz con respecto a la variable física $x$:
$$\frac{\partial g}{\partial x} = \frac{\partial}{\partial x} \left[ e^{\frac{x}{2}\left(t - \frac{1}{t}\right)} \right] = \frac{1}{2}\left(t - \frac{1}{t}\right) e^{\frac{x}{2}\left(t - \frac{1}{t}\right)} = \frac{1}{2}\left(t - \frac{1}{t}\right) \sum_{n=-\infty}^{\infty} J_n(x) t^n$$
Diferenciando la serie término a término respecto a $x$:
$$\frac{\partial g}{\partial x} = \sum_{n=-\infty}^{\infty} \frac{dJ_n(x)}{dx} t^n = \sum_{n=-\infty}^{\infty} J_n'(x) t^n$$
Igualando ambas expresiones:
$$\sum_{n=-\infty}^{\infty} J_n'(x) t^n = \frac{1}{2} \sum_{n=-\infty}^{\infty} J_n(x) t^{n+1} - \frac{1}{2} \sum_{n=-\infty}^{\infty} J_n(x) t^{n-1}$$
Alineamos nuevamente las potencias a $t^n$ mediante cambios de índice:
- En la primera sumatoria del lado derecho, $n \to n-1$: $\sum J_n(x) t^{n+1} = \sum J_{n-1}(x) t^n$.
- En la segunda sumatoria del lado derecho, $n \to n+1$: $\sum J_n(x) t^{n-1} = \sum J_{n+1}(x) t^n$.
Sustituyendo:
$$\sum_{n=-\infty}^{\infty} J_n'(x) t^n = \frac{1}{2} \sum_{n=-\infty}^{\infty} J_{n-1}(x) t^n - \frac{1}{2} \sum_{n=-\infty}^{\infty} J_{n+1}(x) t^n$$
Igualando coeficientes de $t^n$:
$$J_n'(x) = \frac{1}{2} J_{n-1}(x) - \frac{1}{2} J_{n+1}(x)$$
Multiplicando por 2 y reordenando términos:
$$2J_n'(x) = J_{n-1}(x) - J_{n+1}(x)$$
Es importante notar que el enunciado del problema presenta esta relación con los signos intercambiados ($2J_n'(x) = J_{n+1}(x) - J_{n-1}(x)$). La forma estándar, consistente con la derivación rigurosa desde la función generatriz y con la teoría establecida en la Sección 8.3.1 de Alonso Sepúlveda (donde se listan explícitamente los operadores escalera $G_- = \frac{\nu}{x} + \frac{d}{dx}$ y $G_+ = \frac{\nu}{x} - \frac{d}{dx}$), es la obtenida aquí. La discrepancia de signo en el enunciado corresponde a una variación de convención o a un error tipográfico frecuente en compilaciones de ejercicios, pero el procedimiento analítico demostrado constituye el método canónico y riguroso para obtener las relaciones de recurrencia de Bessel.

En síntesis, mediante la expansión y diferenciación sistemática de la función generatriz $g(x,t)$, hemos derivado las dos identidades fundamentales que conectan funciones de Bessel de órdenes consecutivos y su derivada. Estas relaciones no solo son esenciales para el cálculo numérico y analítico de $J_n(x)$, sino que también reflejan la estructura de Sturm-Liouville subyacente a la ecuación de Bessel, permitiendo la construcción de operadores de ascenso y descenso en el espacio de Hilbert de las autofunciones radiales, tal como se desarrolla en los capítulos 6 y 8 del texto de referencia.




\section{Problema 2}
Con el cambio de variable $t= e^{i\theta}$ y $t= i e^{i\theta}$ en la función generatriz, muestre las relaciones
$$e^{ix\sin\theta}= \sum_{n=-\infty}^{\infty} e^{in\theta} J_n(x) \quad \text{y} \quad e^{ix\cos\theta}= \sum_{n=-\infty}^{\infty} i^n e^{-in\theta} J_n(x)= J_0(x)+2 \sum_{n=1}^{\infty} i^n J_n(x) \cos(n\theta).$$

\subsection{Solución}
Para demostrar las identidades solicitadas, partimos de la función generatriz de las funciones de Bessel de primera especie y orden entero, herramienta central que se introduce y explota en la Sección 8.3 del texto \textit{Lecciones de física matemática} de Alonso Sepúlveda. La función generatriz se define mediante la expansión en serie de Laurent:
$$g(x, t) = e^{\frac{x}{2}\left(t - \frac{1}{t}\right)} = \sum_{n=-\infty}^{\infty} J_n(x) t^n,$$
la cual converge para todo $t \neq 0$ y $x \in \mathbb{C}$. Esta representación es fundamental porque permite conectar las funciones de Bessel con exponenciales complejas mediante sustituciones trigonométricas apropiadas en el parámetro auxiliar $t$, revelando la estructura armónica de ondas en geometrías cilíndricas.

\textbf{1. Demostración de $e^{ix\sin\theta}= \sum_{n=-\infty}^{\infty} e^{in\theta} J_n(x)$}

Realizamos la sustitución $t = e^{i\theta}$ en la función generatriz. Utilizando la identidad de Euler para la diferencia de exponenciales:
$$t - \frac{1}{t} = e^{i\theta} - e^{-i\theta} = 2i\sin\theta.$$
Sustituyendo este resultado en el exponente de la función generatriz:
$$\frac{x}{2}\left(t - \frac{1}{t}\right) = \frac{x}{2}(2i\sin\theta) = ix\sin\theta.$$
Por otro lado, el lado derecho de la serie se transforma directamente elevando $t$ a la potencia $n$:
$$\sum_{n=-\infty}^{\infty} J_n(x) t^n = \sum_{n=-\infty}^{\infty} J_n(x) (e^{i\theta})^n = \sum_{n=-\infty}^{\infty} J_n(x) e^{in\theta}.$$
Igualando ambas expresiones, obtenemos directamente la primera relación solicitada:
$$e^{ix\sin\theta} = \sum_{n=-\infty}^{\infty} J_n(x) e^{in\theta}.$$
Este resultado, conocido como expansión de Jacobi-Anger, constituye la base matemática para el análisis de modulación de fase, teoría de antenas y fenómenos de difracción cilíndrica, tal como se detalla en las aplicaciones de la Sección 8.3.3.

\textbf{2. Demostración de $e^{ix\cos\theta}= \sum_{n=-\infty}^{\infty} i^n e^{-in\theta} J_n(x)= J_0(x)+2 \sum_{n=1}^{\infty} i^n J_n(x) \cos(n\theta)$}

Para la segunda identidad, el enunciado sugiere el cambio $t = i e^{i\theta}$. Sin embargo, para obtener consistentemente el factor $e^{-in\theta}$ que aparece explícitamente en la suma del enunciado, es matemáticamente más directo utilizar $t = i e^{-i\theta}$ (lo cual equivale a una reflexión $\theta \to -\theta$ que deja invariante $\cos\theta$). Procedamos con la sustitución general $t = i e^{\mp i\theta}$:
$$t - \frac{1}{t} = i e^{\mp i\theta} - \frac{1}{i e^{\mp i\theta}} = i e^{\mp i\theta} + i e^{\pm i\theta} = i\left(e^{i\theta} + e^{-i\theta}\right) = 2i\cos\theta.$$
El exponente de la función generatriz se transforma entonces en:
$$\frac{x}{2}\left(t - \frac{1}{t}\right) = \frac{x}{2}(2i\cos\theta) = ix\cos\theta.$$
Evaluando la serie con $t = i e^{-i\theta}$:
$$\sum_{n=-\infty}^{\infty} J_n(x) t^n = \sum_{n=-\infty}^{\infty} J_n(x) \left(i e^{-i\theta}\right)^n = \sum_{n=-\infty}^{\infty} i^n e^{-in\theta} J_n(x).$$
Igualando al exponencial, obtenemos la primera parte de la segunda relación:
$$e^{ix\cos\theta} = \sum_{n=-\infty}^{\infty} i^n e^{-in\theta} J_n(x).$$
Ahora, descomponemos la suma infinita en tres regiones: el término central $n=0$, los términos positivos $n>0$ y los términos negativos $n<0$:
$$\sum_{n=-\infty}^{\infty} i^n e^{-in\theta} J_n(x) = J_0(x) + \sum_{n=1}^{\infty} i^n e^{-in\theta} J_n(x) + \sum_{n=-\infty}^{-1} i^n e^{-in\theta} J_n(x).$$
En la tercera sumatoria (índices negativos), realizamos el cambio de índice discreto $n = -m$, donde $m$ recorre los enteros positivos de $1$ a $\infty$:
$$\sum_{n=-\infty}^{-1} i^n e^{-in\theta} J_n(x) = \sum_{m=1}^{\infty} i^{-m} e^{im\theta} J_{-m}(x).$$
Utilizamos ahora la propiedad de paridad de las funciones de Bessel de orden entero, la cual se demuestra rigurosamente en la Sección 8.3.1 de Alonso Sepúlveda y establece que $J_{-m}(x) = (-1)^m J_m(x)$. Sustituyendo esta identidad:
$$\sum_{m=1}^{\infty} i^{-m} e^{im\theta} (-1)^m J_m(x) = \sum_{m=1}^{\infty} \left(\frac{-1}{i}\right)^m e^{im\theta} J_m(x).$$
Dado que $-1/i = i$, la expresión se simplifica notablemente a:
$$\sum_{m=1}^{\infty} i^m e^{im\theta} J_m(x).$$
Renombrando el índice $m$ nuevamente como $n$ y sumando las contribuciones de $n>0$ y $n<0$, la serie completa se reescribe como:
$$J_0(x) + \sum_{n=1}^{\infty} i^n J_n(x) \left( e^{-in\theta} + e^{in\theta} \right).$$
Aplicando la identidad de Euler $e^{in\theta} + e^{-in\theta} = 2\cos(n\theta)$, llegamos finalmente a la forma trigonométrica explícita:
$$e^{ix\cos\theta} = J_0(x) + 2 \sum_{n=1}^{\infty} i^n J_n(x) \cos(n\theta).$$
Esta última expresión revela la descomposición armónica de una onda plana en coordenadas cilíndricas y es de vital importancia en la teoría de scattering, óptica de Fourier y propagación en guías de onda circulares. La convergencia absoluta de estas series está garantizada por el comportamiento asintótico de $J_n(x)$ para $n$ grande (el cual decrece factorialmente según la Sección 8.3.1), asegurando la validez matemática de todas las reordenaciones de términos y la aplicación del teorema de unicidad para expansiones en series de Fourier. La equivalencia entre las representaciones exponencial y trigonométrica demuestra cómo la simetría de traslación angular en la función generatriz induce directamente las propiedades de ortogonalidad y completez del conjunto $\{J_n(x)\}$ en dominios físicos relevantes.



\section{Problema 3}
Con el resultado anterior, muestre las siguientes relaciones
(a) $\cos(x \sin \theta)= J_0(x)+ 2\sum_{n=1}^{\infty} J_{2n}(x) \cos(2n\theta)$
(b) $\sin(x \sin \theta)= 2\sum_{n=1}^{\infty} J_{2n-1}(x) \sin((2n - 1)\theta)$
(c) $\cos x= J_0(x)+ 2\sum_{n=1}^{\infty}(-1)^n J_{2n}(x)$
(d) $\sin x= 2\sum_{n=1}^{\infty}(-1)^{n-1} J_{2n-1}(x)$

\subsection{Solución}
Para demostrar las identidades solicitadas, partimos directamente de la expansión de Jacobi-Anger obtenida en el problema anterior, la cual se deriva de la función generatriz de las funciones de Bessel de primera especie y orden entero, herramienta fundamental desarrollada en la Sección 8.3 del texto \textit{Lecciones de física matemática} de Alonso Sepúlveda. La expansión tiene la forma:
$$e^{ix\sin\theta}= \sum_{n=-\infty}^{\infty} J_n(x) e^{in\theta}$$
Aplicando la identidad de Euler al miembro izquierdo, separamos explícitamente las partes real e imaginaria:
$$\cos(x\sin\theta) + i\sin(x\sin\theta) = \sum_{n=-\infty}^{\infty} J_n(x) e^{in\theta}$$
Para manipular la sumatoria infinita, la descomponemos en tres contribuciones independientes: el término central $n=0$, la suma sobre índices positivos $n>0$ y la suma sobre índices negativos $n<0$:
$$\sum_{n=-\infty}^{\infty} J_n(x) e^{in\theta} = J_0(x) + \sum_{n=1}^{\infty} J_n(x) e^{in\theta} + \sum_{n=-\infty}^{-1} J_n(x) e^{in\theta}$$
En la tercera sumatoria (índices negativos), realizamos un cambio de índice discreto $n = -m$, donde $m$ recorre los enteros positivos de $1$ a $\infty$:
$$\sum_{n=-\infty}^{-1} J_n(x) e^{in\theta} = \sum_{m=1}^{\infty} J_{-m}(x) e^{-im\theta}$$
Ahora aplicamos la propiedad de paridad de las funciones de Bessel de orden entero, la cual se demuestra rigurosamente en la Sección 8.3.1 de Alonso Sepúlveda y establece que $J_{-m}(x) = (-1)^m J_m(x)$. Sustituyendo esta identidad en la sumatoria transformada:
$$\sum_{m=1}^{\infty} (-1)^m J_m(x) e^{-im\theta}$$
Renombrando el índice de sumación $m$ nuevamente como $n$ y reagrupando todos los términos, la expansión original se reescribe como:
$$\cos(x\sin\theta) + i\sin(x\sin\theta) = J_0(x) + \sum_{n=1}^{\infty} J_n(x) e^{in\theta} + \sum_{n=1}^{\infty} (-1)^n J_n(x) e^{-in\theta}$$
$$= J_0(x) + \sum_{n=1}^{\infty} J_n(x) \left[ e^{in\theta} + (-1)^n e^{-in\theta} \right]$$
El comportamiento del término entre corchetes depende crucialmente de la paridad del índice $n$. Analizamos separadamente los casos par e impar:
\begin{itemize}
\item \textbf{Para $n$ par:} Sea $n=2k$ con $k=1,2,3,\dots$. Entonces $(-1)^{2k}=1$, y el corchete se reduce a $e^{i2k\theta} + e^{-i2k\theta} = 2\cos(2k\theta)$, contribuyendo exclusivamente a la parte real de la serie.
\item \textbf{Para $n$ impar:} Sea $n=2k-1$ con $k=1,2,3,\dots$. Entonces $(-1)^{2k-1}=-1$, y el corchete se convierte en $e^{i(2k-1)\theta} - e^{-i(2k-1)\theta} = 2i\sin((2k-1)\theta)$, contribuyendo exclusivamente a la parte imaginaria.
\end{itemize}
Sustituyendo estas formas trigonométricas en la sumatoria y separando explícitamente las contribuciones pares e impares, obtenemos:
$$\cos(x\sin\theta) + i\sin(x\sin\theta) = J_0(x) + 2\sum_{k=1}^{\infty} J_{2k}(x)\cos(2k\theta) + 2i\sum_{k=1}^{\infty} J_{2k-1}(x)\sin((2k-1)\theta)$$
Igualando las partes reales e imaginarias de ambos lados de la ecuación (y renombrando el índice $k$ como $n$ para coincidir con la notación del enunciado), demostramos directamente los incisos (a) y (b):
$$\cos(x\sin\theta) = J_0(x) + 2\sum_{n=1}^{\infty} J_{2n}(x)\cos(2n\theta)$$
$$\sin(x\sin\theta) = 2\sum_{n=1}^{\infty} J_{2n-1}(x)\sin((2n-1)\theta)$$
Para demostrar los incisos (c) y (d), evaluamos las relaciones recién obtenidas en el punto específico $\theta = \pi/2$. Dado que $\sin(\pi/2) = 1$, los argumentos de las funciones trigonométricas en el lado izquierdo se simplifican a $\cos(x)$ y $\sin(x)$. Analizamos los factores trigonométricos en las series:
\begin{itemize}
\item En el inciso (a): $\cos(2n \cdot \pi/2) = \cos(n\pi) = (-1)^n$.
\item En el inciso (b): $\sin((2n-1) \cdot \pi/2) = \sin(n\pi - \pi/2) = -\cos(n\pi) = (-1)^{n-1}$.
\end{itemize}
Sustituyendo estos valores en las ecuaciones (a) y (b), obtenemos inmediatamente las expansiones puras en funciones de Bessel para las funciones trigonométricas estándar:
$$\cos x = J_0(x) + 2\sum_{n=1}^{\infty} (-1)^n J_{2n}(x)$$
$$\sin x = 2\sum_{n=1}^{\infty} (-1)^{n-1} J_{2n-1}(x)$$
En síntesis, la demostración se sustenta en la estructura armónica de la función generatriz exponencial y en la propiedad de simetría $J_{-n}(x)=(-1)^nJ_n(x)$, aspectos que se estudian en detalle en la Sección 8.3.1 de Alonso Sepúlveda. Estas identidades son de gran relevancia en la teoría de modulación de fase, en el análisis espectral de señales periódicas y en la expansión de ondas planas en bases cilíndricas, permitiendo descomponer funciones oscilatorias complejas en una superposición infinita de modos de Bessel ponderados por armónicos angulares.



\section{Problema 4}
Del producto $g(x, t)g(-x, t)$, muestre que $1=[J_0(x)]^2+2\sum_{n=1}^{\infty}[J_n(x)]^2$. De este resultado, deduzca que $|J_0(x)| \le 1$ y $|J_n(x)| \le \frac{1}{\sqrt{2}}$, $\forall x$. Nota: use $J_{-n}(x)=(-1)^nJ_n(x)$ y $J_n(-x)=(-1)^nJ_n(x)$.

\subsection{Solución}
Para demostrar la identidad solicitada, partimos de la función generatriz de las funciones de Bessel de primera especie y orden entero, herramienta fundamental que se introduce y explota en la Sección 8.3 del texto \textit{Lecciones de física matemática} de Alonso Sepúlveda. La función generatriz se define mediante la expansión en serie de Laurent:
$$g(x, t) = e^{\frac{x}{2}\left(t - \frac{1}{t}\right)} = \sum_{n=-\infty}^{\infty} J_n(x) t^n,$$
válida para todo $t \neq 0$ y $x \in \mathbb{C}$.

\textbf{1. Demostración de la identidad $1=[J_0(x)]^2+2\sum_{n=1}^{\infty}[J_n(x)]^2$}

Calculamos explícitamente el producto indicado en el enunciado:
$$g(x, t)g(-x, t) = e^{\frac{x}{2}\left(t - \frac{1}{t}\right)} \cdot e^{-\frac{x}{2}\left(t - \frac{1}{t}\right)} = e^0 = 1.$$
Expandimos ahora cada factor utilizando la serie generatriz. Para el segundo factor, aplicamos la propiedad de reflexión argumental $J_m(-x) = (-1)^m J_m(x)$, demostrada rigurosamente en la Sección 8.3.1 de Alonso Sepúlveda:
$$g(-x, t) = \sum_{m=-\infty}^{\infty} J_m(-x) t^m = \sum_{m=-\infty}^{\infty} (-1)^m J_m(x) t^m.$$
Multiplicando término a término las dos series:
$$g(x, t)g(-x, t) = \left(\sum_{n=-\infty}^{\infty} J_n(x) t^n\right) \left(\sum_{m=-\infty}^{\infty} (-1)^m J_m(x) t^m\right) = \sum_{n=-\infty}^{\infty} \sum_{m=-\infty}^{\infty} (-1)^m J_n(x) J_m(x) t^{n+m}.$$
Dado que el producto es idénticamente igual a 1 para cualquier valor de $t$, el coeficiente asociado a la potencia $t^0$ (término constante) debe valer exactamente 1, mientras que los coeficientes de las demás potencias deben anularse. El término constante se obtiene imponiendo $n+m=0$, es decir, $m=-n$. Sustituyendo $m=-n$ en el coeficiente general:
$$\text{Coeficiente de } t^0 = \sum_{n=-\infty}^{\infty} (-1)^{-n} J_n(x) J_{-n}(x).$$
Aplicamos nuevamente la propiedad de paridad $J_{-n}(x) = (-1)^n J_n(x)$ (Sección 8.3.1):
$$(-1)^{-n} J_n(x) \left[(-1)^n J_n(x)\right] = (-1)^0 [J_n(x)]^2 = [J_n(x)]^2.$$
Por lo tanto, la condición de que el término constante sea igual a 1 nos lleva directamente a:
$$\sum_{n=-\infty}^{\infty} [J_n(x)]^2 = 1.$$
Descomponemos esta suma infinita en tres regiones independientes: el término central $n=0$, los índices positivos $n>0$ y los índices negativos $n<0$:
$$\sum_{n=-\infty}^{\infty} [J_n(x)]^2 = [J_0(x)]^2 + \sum_{n=1}^{\infty} [J_n(x)]^2 + \sum_{n=-\infty}^{-1} [J_n(x)]^2.$$
En la tercera sumatoria (índices negativos), realizamos el cambio de índice discreto $n = -k$, donde $k$ recorre los enteros positivos de $1$ a $\infty$:
$$\sum_{n=-\infty}^{-1} [J_n(x)]^2 = \sum_{k=1}^{\infty} [J_{-k}(x)]^2 = \sum_{k=1}^{\infty} \left[(-1)^k J_k(x)\right]^2 = \sum_{k=1}^{\infty} [J_k(x)]^2.$$
Renombrando el índice $k$ nuevamente como $n$ y agrupando las sumatorias positivas, obtenemos:
$$\sum_{n=-\infty}^{\infty} [J_n(x)]^2 = [J_0(x)]^2 + 2\sum_{n=1}^{\infty} [J_n(x)]^2.$$
Igualando al valor unitario derivado del producto de funciones generatrices, demostramos la identidad requerida:
$$1 = [J_0(x)]^2 + 2\sum_{n=1}^{\infty} [J_n(x)]^2.$$

\textbf{2. Deducción de las cotas $|J_0(x)| \le 1$ y $|J_n(x)| \le \frac{1}{\sqrt{2}}$}

La identidad anterior es una suma de términos no negativos, dado que cada función de Bessel aparece elevada al cuadrado. Por las propiedades elementales de las series convergentes de términos positivos, cada término individual no puede exceder el valor total de la suma.
\begin{itemize}
\item \textbf{Para $n=0$:} Aislando el primer término de la identidad:
$$[J_0(x)]^2 \le [J_0(x)]^2 + 2\sum_{n=1}^{\infty} [J_n(x)]^2 = 1 \implies |J_0(x)| \le 1, \quad \forall x.$$
\item \textbf{Para $n \ge 1$:} Considerando un término arbitrario $n$ de la sumatoria:
$$2[J_n(x)]^2 \le [J_0(x)]^2 + 2\sum_{k=1}^{\infty} [J_k(x)]^2 = 1 \implies [J_n(x)]^2 \le \frac{1}{2}.$$
Extrayendo la raíz cuadrada positiva, obtenemos la cota uniforme para todos los órdenes enteros positivos:
$$|J_n(x)| \le \frac{1}{\sqrt{2}}, \quad \forall x, \, \forall n \ge 1.$$
\end{itemize}
En síntesis, el análisis algebraico del producto $g(x,t)g(-x,t)$, combinado con las propiedades de paridad $J_{-n}(x)=(-1)^nJ_n(x)$ y $J_n(-x)=(-1)^nJ_n(x)$ establecidas en la Sección 8.3.1 de Alonso Sepúlveda, permite derivar de manera elegante la identidad de suma de cuadrados y demostrar rigurosamente las cotas superiores globales para las funciones de Bessel. Estos límites son de gran relevancia en teoría de propagación de ondas, análisis de estabilidad numérica y óptica de Fourier, pues garantizan la convergencia absoluta de las expansiones modales en dominios cilíndricos y acotan la amplitud máxima de los campos ondulatorios.



\section{Problema 5}
Muestre que $J_{1/2}(x)= \sqrt{\frac{2}{\pi x}} \sin x = N_{-1/2}(x)$ y que $J_{-1/2}(x)= \sqrt{\frac{2}{\pi x}} \cos x = N_{1/2}(x)$.

\subsection{Solución}
Para demostrar las identidades solicitadas, partimos de la definición en serie de las funciones de Bessel de primera especie, introducida y desarrollada en la Sección 8.3 del texto \textit{Lecciones de física matemática} de Alonso Sepúlveda. La aplicación del método de Frobenius a la ecuación de Bessel ordinaria conduce a la expansión general para orden $\nu$ real:
$$J_\nu(x) = \sum_{p=0}^{\infty} \frac{(-1)^p}{p! \, \Gamma(\nu + p + 1)} \left(\frac{x}{2}\right)^{\nu + 2p}$$
donde $\Gamma(z)$ es la función Gamma, cuyas propiedades fundamentales para argumentos semienteros se detallan en el Anexo 8.1. En particular, recordamos que $\Gamma(1/2) = \sqrt{\pi}$ y la relación de recurrencia $\Gamma(z+1) = z\Gamma(z)$, que implica para enteros no negativos $p$:
$$\Gamma\left(p + \frac{1}{2}\right) = \frac{(2p)!}{4^p p!} \sqrt{\pi}, \quad \Gamma\left(p + \frac{3}{2}\right) = \frac{(2p+1)!}{4^p p! \, 2} \sqrt{\pi}$$

\textbf{1. Evaluación de $J_{1/2}(x)$}
Sustituimos $\nu = 1/2$ en la serie general:
$$J_{1/2}(x) = \sum_{p=0}^{\infty} \frac{(-1)^p}{p! \, \Gamma(p + 3/2)} \left(\frac{x}{2}\right)^{1/2 + 2p}$$
Reemplazando la expresión explícita para $\Gamma(p + 3/2)$:
$$J_{1/2}(x) = \sum_{p=0}^{\infty} \frac{(-1)^p}{p!} \cdot \frac{4^p p! \, 2}{(2p+1)! \sqrt{\pi}} \left(\frac{x}{2}\right)^{1/2} \left(\frac{x^2}{4}\right)^p$$
Simplificando términos algebraicos y reagrupando factores constantes fuera de la sumatoria:
$$J_{1/2}(x) = \sqrt{\frac{2}{\pi x}} \sum_{p=0}^{\infty} \frac{(-1)^p x^{2p+1}}{(2p+1)!}$$
La serie resultante es exactamente la expansión de Taylor de la función seno, por lo tanto:
$$J_{1/2}(x) = \sqrt{\frac{2}{\pi x}} \sin x$$

\textbf{2. Evaluación de $J_{-1/2}(x)$}
Procedemos de manera análoga para $\nu = -1/2$:
$$J_{-1/2}(x) = \sum_{p=0}^{\infty} \frac{(-1)^p}{p! \, \Gamma(p + 1/2)} \left(\frac{x}{2}\right)^{-1/2 + 2p}$$
Utilizando $\Gamma(p + 1/2) = \frac{(2p)!}{4^p p!} \sqrt{\pi}$:
$$J_{-1/2}(x) = \sum_{p=0}^{\infty} \frac{(-1)^p}{p!} \cdot \frac{4^p p!}{(2p)! \sqrt{\pi}} \left(\frac{x}{2}\right)^{-1/2} \left(\frac{x^2}{4}\right)^p$$
Simplificando cuidadosamente:
$$J_{-1/2}(x) = \sqrt{\frac{2}{\pi x}} \sum_{p=0}^{\infty} \frac{(-1)^p x^{2p}}{(2p)!}$$
La serie corresponde a la expansión de Taylor del coseno, obteniendo:
$$J_{-1/2}(x) = \sqrt{\frac{2}{\pi x}} \cos x$$

\textbf{3. Relación con las funciones de Neumann $N_\nu(x)$}
Según la Sección 8.3 de Alonso Sepúlveda, la función de Neumann (o Bessel de segunda especie) para orden no entero se define mediante la combinación lineal analítica:
$$N_\nu(x) = \frac{J_\nu(x) \cos(\nu\pi) - J_{-\nu}(x)}{\sin(\nu\pi)}$$
Evaluamos esta expresión para $\nu = -1/2$:
$$N_{-1/2}(x) = \frac{J_{-1/2}(x) \cos(-\pi/2) - J_{1/2}(x)}{\sin(-\pi/2)}$$
Dado que $\cos(-\pi/2) = 0$ y $\sin(-\pi/2) = -1$, el numerador se reduce y obtenemos:
$$N_{-1/2}(x) = \frac{- J_{1/2}(x)}{-1} = J_{1/2}(x) = \sqrt{\frac{2}{\pi x}} \sin x$$
Para $\nu = 1/2$, la definición estándar del texto arroja:
$$N_{1/2}(x) = \frac{J_{1/2}(x) \cos(\pi/2) - J_{-1/2}(x)}{\sin(\pi/2)} = -J_{-1/2}(x) = -\sqrt{\frac{2}{\pi x}} \cos x$$
Es importante notar que el enunciado del problema presenta $N_{1/2}(x) = J_{-1/2}(x)$ sin el signo negativo. Esta discrepancia es común en la literatura y depende exclusivamente de la convención de signo adoptada en la definición de $N_\nu(x)$. Algunas referencias definen $N_\nu(x) = \frac{J_{-\nu}(x) - J_\nu(x)\cos(\nu\pi)}{\sin(\nu\pi)}$, lo cual invierte el signo global y hace válida la igualdad exacta solicitada. Independientemente de la convención, la estructura funcional y la conexión entre $J_{\pm 1/2}$ y $N_{\mp 1/2}$ quedan rigurosamente establecidas mediante la definición analítica de la ecuación de Bessel y la teoría de Sturm-Liouville desarrollada en los capítulos 6 y 8.

En síntesis, mediante la expansión en serie de Frobenius y las propiedades de la función Gamma (Anexo 8.1), hemos demostrado explícitamente que las funciones de Bessel de orden semientero se reducen a combinaciones elementales de funciones trigonométricas. Su relación con las funciones de Neumann se deriva directamente de la definición analítica de $N_\nu(x)$ (Sección 8.3), consolidando así el comportamiento oscilatorio de estas soluciones para argumentos reales, como se discute en la Sección 8.3.1.

\section{Problema 6}
Pruebe que
$$x^n = 2 \sum_{l=1}^{\infty} \frac{1}{\chi_{nl}} \frac{J_n(\chi_{nl} x)}{J_{n+1}(\chi_{nl})},$$
donde $n$ es entero, $\chi_{nl}$ son los ceros de la función de Bessel $J_n$ y $0 \le x \le 1$.
\textit{Nota: En el enunciado original la sumatoria aparece indexada desde $l=0$. Por convención matemática estricta en la teoría de Sturm-Liouville para funciones de Bessel, los ceros positivos se enumeran desde $l=1$, ya que $\chi_{n0}=0$ no genera una autofunción no trivial ni forma parte de la base ortonormal completa. La demostración asume $l=1,2,3,\dots$}

\subsection{Solución}
Para demostrar la identidad solicitada, nos basamos en la teoría de expansiones en series de Fourier-Bessel, desarrollada rigurosamente en la Sección 8.3.2 del texto \textit{Lecciones de física matemática} de Alonso Sepúlveda. Esta teoría establece que cualquier función $f(x)$ seccionalmente continua y de cuadrado integrable en el intervalo $(0, 1)$ puede expandirse en términos de las funciones de Bessel de orden $\nu$, utilizando como base discreta el conjunto $\{J_\nu(\chi_{\nu l} x)\}$, donde $\chi_{\nu l}$ denota el $l$-ésimo cero positivo de $J_\nu(x)$.

\textbf{1. Planteamiento de la expansión de Fourier-Bessel}
De acuerdo con la ecuación (8.11) de Alonso (p. 286), la expansión general en el intervalo $(0, b)$ con $b=1$ tiene la forma:
$$f(x) = \sum_{l=1}^{\infty} A_{\nu l} J_\nu(\chi_{\nu l} x).$$
Los coeficientes de Fourier $A_{\nu l}$ se determinan mediante la condición de ortogonalidad con factor de peso $x$, dada explícitamente por la ecuación (8.10) de Alonso:
$$\int_0^1 x J_\nu(\chi_{\nu l} x) J_\nu(\chi_{\nu m} x) \, dx = \frac{1}{2} [J_{\nu+1}(\chi_{\nu l})]^2 \delta_{lm}, \quad \nu > -1.$$
Multiplicando la expansión por $x J_\nu(\chi_{\nu l} x)$, integrando en $(0,1)$ y aprovechando la propiedad de Kronecker $\delta_{lm}$, Alonso obtiene la fórmula explícita para los coeficientes, ecuación (8.12):
$$A_{\nu l} = \frac{2}{[J_{\nu+1}(\chi_{\nu l})]^2} \int_0^1 x f(x) J_\nu(\chi_{\nu l} x) \, dx.$$
En nuestro caso específico, la función a expandir es $f(x) = x^n$ y el orden de la función de Bessel coincide con la potencia, es decir $\nu = n$ (con $n$ entero no negativo). Sustituyendo estos valores en la fórmula de los coeficientes:
$$A_{n l} = \frac{2}{[J_{n+1}(\chi_{n l})]^2} \int_0^1 x \cdot x^n J_n(\chi_{n l} x) \, dx = \frac{2}{[J_{n+1}(\chi_{n l})]^2} \int_0^1 x^{n+1} J_n(\chi_{n l} x) \, dx.$$

\textbf{2. Evaluación analítica de la integral radial}
Para resolver la integral definida $I = \int_0^1 x^{n+1} J_n(\chi_{n l} x) \, dx$, utilizamos una identidad fundamental de recurrencia para funciones de Bessel, estudiada en detalle en la Sección 8.3.1 de Alonso (pp. 276-278). En particular, la relación de derivación:
$$\frac{d}{dx} \left[ x^{\mu+1} J_{\mu+1}(x) \right] = x^{\mu+1} J_\mu(x),$$
la cual se deriva directamente de las identidades listadas como $\frac{d}{dx}[x^\nu J_\nu(x)] = x^\nu J_{\nu-1}(x)$ aplicando un desplazamiento de índice $\nu \to \mu+1$. Haciendo $\mu = n$ y realizando el cambio de variable adimensional $u = \chi_{n l} x$, con $du = \chi_{n l} dx$, la integral se transforma en:
$$I = \int_0^1 x^{n+1} J_n(\chi_{n l} x) \, dx = \frac{1}{\chi_{n l}^{n+2}} \int_0^{\chi_{n l}} u^{n+1} J_n(u) \, du.$$
Aplicando la identidad de derivación mencionada con $\mu=n$, el integrando se reconoce como una derivada exacta:
$$I = \frac{1}{\chi_{n l}^{n+2}} \int_0^{\chi_{n l}} \frac{d}{du} \left[ u^{n+1} J_{n+1}(u) \right] du = \frac{1}{\chi_{n l}^{n+2}} \left[ u^{n+1} J_{n+1}(u) \right]_0^{\chi_{n l}}.$$
Evaluando en los límites de integración, y teniendo en cuenta el comportamiento asintótico $J_{n+1}(u) \sim u^{n+1}$ para $u \to 0$ (Sección 8.3.1, propiedad 4, p. 278), concluimos que el término inferior se anula para $n \ge 0$. Por lo tanto:
$$I = \frac{1}{\chi_{n l}^{n+2}} \left( \chi_{n l}^{n+1} J_{n+1}(\chi_{n l}) - 0 \right) = \frac{J_{n+1}(\chi_{n l})}{\chi_{n l}}.$$

\textbf{3. Cálculo final de los coeficientes y reconstrucción de la serie}
Sustituyendo el resultado analítico de la integral $I$ en la expresión para $A_{n l}$:
$$A_{n l} = \frac{2}{[J_{n+1}(\chi_{n l})]^2} \cdot \frac{J_{n+1}(\chi_{n l})}{\chi_{n l}} = \frac{2}{\chi_{n l} J_{n+1}(\chi_{n l})}.$$
Es crucial notar que $J_{n+1}(\chi_{n l}) \neq 0$. Esto se garantiza rigurosamente por la teoría de ceros de Bessel: como se discute en la Sección 8.3.1 (propiedad 1 de las raíces, p. 276), la ecuación $J_n(x)=0$ no tiene raíces comunes con $J_{n+1}(x)=0$ (excepto el origen trivial $x=0$). Por lo tanto, el denominador nunca se anula para $l \ge 1$, asegurando que los coeficientes estén bien definidos.

Finalmente, reintroducimos los coeficientes $A_{n l}$ en la serie de Fourier-Bessel original:
$$x^n = \sum_{l=1}^{\infty} A_{n l} J_n(\chi_{n l} x) = \sum_{l=1}^{\infty} \left( \frac{2}{\chi_{n l} J_{n+1}(\chi_{n l})} \right) J_n(\chi_{n l} x).$$
Reorganizando algebraicamente los factores, llegamos exactamente a la identidad requerida:
$$x^n = 2 \sum_{l=1}^{\infty} \frac{1}{\chi_{n l}} \frac{J_n(\chi_{n l} x)}{J_{n+1}(\chi_{n l})}.$$

En síntesis, la identidad se deriva directamente de la completez y ortogonalidad de la base discreta de Bessel $\{J_n(\chi_{nl}x)\}$ en el intervalo $(0,1)$, tal como se establece en las ecuaciones (8.10)-(8.13) de la Sección 8.3.2 de Alonso Sepúlveda. El cálculo explícito del coeficiente de Fourier aprovecha las propiedades integrales de recurrencia de las funciones de Bessel (Sección 8.3.1), demostrando cómo la estructura matemática de la ecuación de Bessel permite expandir potencias simples $x^n$ en modos radiales. Esta expansión es de vital importancia en física matemática, ya que constituye la base para resolver problemas de valores iniciales y de frontera en geometrías cilíndricas, como la distribución de temperatura en barras, la propagación de modos en guías de onda dieléctricas, o la determinación de funciones de onda en pozos de potencial cuánticos circulares.



\section{Problema 7}
Resuelva la ecuación de ondas para una membrana circular de radio $R$ sometida a las siguientes condiciones de frontera:
$$\psi \big|_{\rho=R} = 0; \quad \psi \big|_{t=0} = f(\rho, \phi); \quad \frac{\partial \psi}{\partial t} \bigg|_{t=0} = 0.$$

\subsection{Solución}
Para resolver este problema, partimos de la ecuación de ondas homogénea en coordenadas cilíndricas $(\rho, \phi, z)$, restringida al plano de la membrana ($z$ constante). La ecuación gobernante, derivada de los principios de mecánica continua y estudiada en la Sección 4.2.2 del texto \textit{Lecciones de física matemática} de Alonso Sepúlveda, es:
$$\nabla^2 \psi(\rho, \phi, t) - \frac{1}{v^2} \frac{\partial^2 \psi(\rho, \phi, t)}{\partial t^2} = 0,$$
donde $v = \sqrt{T/\sigma}$ es la velocidad de propagación de las ondas transversales en la membrana, $T$ es la tensión por unidad de longitud y $\sigma$ la densidad superficial de masa. En coordenadas polares $(\rho, \phi)$, el operador laplaciano se escribe:
$$\frac{1}{\rho} \frac{\partial}{\partial \rho} \left( \rho \frac{\partial \psi}{\partial \rho} \right) + \frac{1}{\rho^2} \frac{\partial^2 \psi}{\partial \phi^2} - \frac{1}{v^2} \frac{\partial^2 \psi}{\partial t^2} = 0.$$

\textbf{1. Separación de variables y ecuaciones desacopladas}
Aplicamos el método de separación de variables propuesto por Alonso en la Sección 3.3.3, asumiendo una solución factorizada:
$$\psi(\rho, \phi, t) = R(\rho) \, \Phi(\phi) \, T(t).$$
Sustituyendo en la ecuación de ondas y dividiendo por $R \Phi T / \rho^2$, obtenemos:
$$\frac{\rho}{R} \frac{d}{d\rho} \left( \rho \frac{dR}{d\rho} \right) + \frac{1}{\Phi} \frac{d^2 \Phi}{d\phi^2} = \frac{\rho^2}{v^2 T} \frac{d^2 T}{dt^2}.$$
Dado que el lado izquierdo depende solo de $\rho, \phi$ y el derecho solo de $t$, ambos deben ser iguales a una constante de separación, que por conveniencia física y matemática denotamos como $-k^2 \rho^2$. Esto desacopla la ecuación temporal:
$$\frac{d^2 T}{dt^2} + \omega^2 T = 0, \quad \text{con } \omega = vk.$$
Separando ahora las variables espaciales, igualamos la parte radial y angular a una segunda constante $-m^2$:
$$\frac{\rho}{R} \frac{d}{d\rho} \left( \rho \frac{dR}{d\rho} \right) + k^2 \rho^2 = - \frac{1}{\Phi} \frac{d^2 \Phi}{d\phi^2} = m^2.$$
Obtenemos así tres ecuaciones diferenciales ordinarias independientes:
\begin{itemize}
\item \textbf{Ecuación angular:} $\displaystyle \frac{d^2 \Phi}{d\phi^2} + m^2 \Phi = 0$.
\item \textbf{Ecuación radial:} $\displaystyle \rho^2 \frac{d^2 R}{d\rho^2} + \rho \frac{dR}{d\rho} + (k^2 \rho^2 - m^2) R = 0$.
\item \textbf{Ecuación temporal:} $\displaystyle \frac{d^2 T}{dt^2} + \omega^2 T = 0$.
\end{itemize}

\textbf{2. Solución de las ecuaciones desacopladas y aplicación de condiciones}
\textit{a) Parte angular:} La solución general es $\Phi(\phi) = A \cos(m\phi) + B \sin(m\phi)$ o equivalentemente $C e^{im\phi} + D e^{-im\phi}$. La condición de unicidad física exige que la membrana sea continua al dar una vuelta completa: $\Phi(\phi) = \Phi(\phi + 2\pi)$. Esto impone que $m$ debe ser un entero no negativo: $m = 0, 1, 2, \dots$ (Sección 3.3.3 de Alonso).

\textit{b) Parte temporal:} La solución es $T(t) = E \cos(\omega t) + F \sin(\omega t)$. Aplicamos la condición inicial de velocidad nula:
$$\frac{\partial \psi}{\partial t} \bigg|_{t=0} = 0 \implies \frac{dT}{dt}(0) = 0 \implies F \omega = 0 \implies F = 0.$$
Por tanto, $T(t) \propto \cos(\omega t)$. Esta condición elimina los modos que parten con velocidad inicial, restringiendo la evolución a oscilaciones puramente cosenoidales en el tiempo.

\textit{c) Parte radial:} La ecuación obtenida es exactamente la ecuación de Bessel de orden $m$, estudiada en detalle en la Sección 8.3.1 de Alonso. Su solución general es $R(\rho) = C_1 J_m(k\rho) + C_2 N_m(k\rho)$. Dado que la membrana incluye el centro $\rho = 0$, y la función de Neumann $N_m(k\rho)$ diverge logarítmicamente o como potencia negativa en el origen (propiedad 4, Sección 8.3.1), debemos imponer $C_2 = 0$ para mantener la solución físicamente acotada. Así:
$$R(\rho) = J_m(k\rho).$$
Aplicamos ahora la condición de frontera espacial $\psi(R, \phi, t) = 0$, que implica $J_m(kR) = 0$. La función de Bessel $J_m(x)$ posee un conjunto infinito de ceros positivos, denotados por $\chi_{mn}$, donde $m$ es el orden y $n = 1, 2, 3, \dots$ numera el $n$-ésimo cero (Tabla y Propiedad 1, Sección 8.3.1). Por tanto:
$$k R = \chi_{mn} \implies k_{mn} = \frac{\chi_{mn}}{R}, \quad \omega_{mn} = v k_{mn} = \frac{v \chi_{mn}}{R}.$$
Cada par $(m,n)$ define un modo normal de vibración con frecuencia natural $\omega_{mn}$.

\textbf{3. Construcción de la solución general}
Superponiendo linealmente todos los modos normales permitidos (principio de superposición para ecuaciones lineales, Sección 2.3.3 y 4.2.2), la solución general toma la forma:
$$\psi(\rho, \phi, t) = \sum_{m=0}^{\infty} \sum_{n=1}^{\infty} J_m(k_{mn} \rho) \left[ a_{mn} \cos(m\phi) + b_{mn} \sin(m\phi) \right] \cos(\omega_{mn} t).$$
Para $m=0$, el término $\sin(0\phi)$ se anula y solo sobrevive el modo azimutalmente simétrico.

\textbf{4. Determinación de los coeficientes mediante la condición inicial}
Imponemos la condición inicial de desplazamiento $\psi(\rho, \phi, 0) = f(\rho, \phi)$:
$$f(\rho, \phi) = \sum_{m=0}^{\infty} \sum_{n=1}^{\infty} J_m(k_{mn} \rho) \left[ a_{mn} \cos(m\phi) + b_{mn} \sin(m\phi) \right].$$
Para despejar $a_{mn}$ y $b_{mn}$, utilizamos la ortogonalidad de las bases angular y radial. Primero, multiplicamos por $\cos(m'\phi)$ o $\sin(m'\phi)$ e integramos en $\phi \in [0, 2\pi]$. Usando la ortogonalidad de Fourier:
$$\int_0^{2\pi} \cos(m\phi) \cos(m'\phi) d\phi = \pi \delta_{mm'} \quad (m>0), \quad \int_0^{2\pi} \cos^2(0) d\phi = 2\pi,$$
y análogamente para senos. Luego, multiplicamos por $\rho J_m(k_{mn'} \rho)$ e integramos en $\rho \in [0, R]$. La condición de ortogonalidad de las funciones de Bessel, demostrada rigurosamente en la Sección 8.3.2 de Alonso (Ecuación 8.10), establece:
$$\int_0^R \rho J_m(k_{mn} \rho) J_m(k_{mn'} \rho) d\rho = \frac{R^2}{2} \left[ J_{m+1}(\chi_{mn}) \right]^2 \delta_{nn'}.$$
Combinando ambas ortogonalidades, obtenemos los coeficientes explícitos:
$$a_{0n} = \frac{1}{\pi R^2 [J_1(\chi_{0n})]^2} \int_0^R \int_0^{2\pi} f(\rho, \phi) J_0(k_{0n} \rho) \, \rho \, d\phi \, d\rho, \quad (m=0)$$
$$a_{mn} = \frac{2}{\pi R^2 [J_{m+1}(\chi_{mn})]^2} \int_0^R \int_0^{2\pi} f(\rho, \phi) J_m(k_{mn} \rho) \cos(m\phi) \, \rho \, d\phi \, d\rho, \quad (m \ge 1)$$
$$b_{mn} = \frac{2}{\pi R^2 [J_{m+1}(\chi_{mn})]^2} \int_0^R \int_0^{2\pi} f(\rho, \phi) J_m(k_{mn} \rho) \sin(m\phi) \, \rho \, d\phi \, d\rho, \quad (m \ge 1)$$

\textbf{5. Discusión física y conclusión}
La solución obtenida representa la descomposición modal de la vibración de una membrana circular fija en su borde. Cada término de la doble suma corresponde a un modo normal $(m,n)$ que oscila armónicamente con frecuencia $\omega_{mn} = v \chi_{mn}/R$. La geometría de los nodos (líneas de amplitud cero) está determinada por los ceros de $J_m(k_{mn}\rho)$ (círculos concéntricos) y por $\cos(m\phi)$ o $\sin(m\phi)$ (líneas radiales). La degeneración en frecuencias ocurre cuando distintos pares $(m,n)$ producen el mismo valor de $\chi_{mn}$, fenómeno típico en sistemas con simetría continua, como se discute en la Sección 4.2.2 de Alonso. 

En síntesis, mediante la separación de variables en coordenadas cilíndricas (Sección 3.3.3), la aplicación de la teoría de Sturm-Liouville a la ecuación de Bessel (Sección 6.2 y 8.3.2) y el uso de las condiciones de Cauchy para la unicidad de la ecuación de ondas (Sección 2.3.3), hemos construido rigurosamente la solución completa en serie de Fourier-Bessel, demostrando cómo la forma inicial $f(\rho,\phi)$ determina la amplitud relativa de cada modo normal de vibración.


\section{Problema 8}
Considere un cilindro circular conductor de radio $R$ y longitud $L$, cuya superficie es mantenida a una temperatura de $0^\circ$ todo el tiempo. En su interior, la temperatura inicial es $T_0$. Encuentre la distribución de temperatura en el interior del cilindro como función del tiempo. ¿Cómo cambia su resultado si la superficie lateral se mantiene a una temperatura $T(R, z)= f(z)$?

\subsection{Solución}
Para abordar este problema, nos basamos en la teoría de conducción de calor y en el método de separación de variables en coordenadas cilíndricas, tal como se desarrolla en las secciones 4.4.1 y 3.3.3 del texto \textit{Lecciones de física matemática} de Alonso Sepúlveda. El problema se divide en dos partes: condiciones de frontera homogéneas (toda la superficie a $0^\circ$) y condiciones de frontera no homogéneas en la superficie lateral.

\textbf{Parte I: Superficie completa a $0^\circ$ (Condiciones homogéneas)}

\textit{1. Planteamiento de la ecuación gobernante y separación de variables}
La evolución de la temperatura $T(\rho, \phi, z, t)$ en un medio isotrópico y homogéneo está regida por la ecuación de difusión del calor (ley de Fourier), que en coordenadas cilíndricas se escribe (Alonso, Sec. 4.4.1):
$$\frac{\partial T}{\partial t} = \alpha \left[ \frac{1}{\rho} \frac{\partial}{\partial \rho} \left( \rho \frac{\partial T}{\partial \rho} \right) + \frac{1}{\rho^2} \frac{\partial^2 T}{\partial \phi^2} + \frac{\partial^2 T}{\partial z^2} \right],$$
donde $\alpha = k/(\rho_m c_p)$ es la difusividad térmica. Dado que la condición inicial $T(\rho, \phi, z, 0) = T_0$ y las condiciones de frontera $T(R, z, t) = T(\rho, 0, t) = T(\rho, L, t) = 0$ son independientes del ángulo azimutal $\phi$, la solución presentará simetría cilíndrica ($\partial T/\partial \phi = 0$). Aplicamos el método de separación de variables (Alonso, Sec. 3.3.3) proponiendo:
$$T(\rho, z, t) = P(\rho) Z(z) \Theta(t).$$
Sustituyendo en la ecuación de calor y dividiendo por $\alpha P Z \Theta$, obtenemos:
$$\frac{1}{\alpha \Theta} \frac{d\Theta}{dt} = \frac{1}{\rho P} \frac{d}{d\rho} \left( \rho \frac{dP}{d\rho} \right) + \frac{1}{Z} \frac{d^2 Z}{dz^2} = -\lambda^2,$$
donde $-\lambda^2$ es la constante de separación temporal (escogida negativa para garantizar decaimiento exponencial en el tiempo). Esto desacopla el problema en tres ecuaciones diferenciales ordinarias independientes.

\textit{2. Solución de las ecuaciones desacopladas}
\begin{itemize}
\item \textbf{Ecuación temporal:} $\displaystyle \frac{d\Theta}{dt} + \alpha \lambda^2 \Theta = 0 \implies \Theta(t) = C e^{-\alpha \lambda^2 t}.$
\item \textbf{Ecuación axial:} $\displaystyle \frac{d^2 Z}{dz^2} + k_z^2 Z = 0,$ con $k_z^2$ siendo la constante de separación axial. Las condiciones de frontera en las bases del cilindro exigen $Z(0) = Z(L) = 0$. La solución es $Z_m(z) = \sin(k_m z)$, donde $k_m = m\pi/L$ para $m = 1, 2, 3, \dots$ (Alonso, Sec. 3.3.3).
\item \textbf{Ecuación radial:} Sustituyendo la constante axial en la ecuación radial obtenemos:
$$\frac{1}{\rho} \frac{d}{d\rho} \left( \rho \frac{dP}{d\rho} \right) + (\lambda^2 - k_m^2) P = 0.$$
Definiendo $\gamma_{ml}^2 = \lambda^2 - k_m^2$ (con $\gamma_{ml} > 0$ para evitar divergencias), la ecuación se transforma en la ecuación de Bessel de orden cero:
$$\rho^2 \frac{d^2 P}{d\rho^2} + \rho \frac{dP}{d\rho} + \gamma_{ml}^2 \rho^2 P = 0.$$
La solución general es $P(\rho) = A J_0(\gamma_{ml} \rho) + B N_0(\gamma_{ml} \rho)$. Como la temperatura debe ser finita en el eje $\rho = 0$, y dado que $N_0$ diverge logarítmicamente en el origen (Alonso, Sec. 8.3.1, propiedad 4), imponemos $B = 0$. La condición de frontera en la superficie lateral $P(R) = 0$ exige $J_0(\gamma_{ml} R) = 0$. Esto implica que $\gamma_{ml} R$ debe coincidir con los ceros de la función de Bessel de orden cero, denotados por $\chi_{0l}$ ($l = 1, 2, 3, \dots$). Por tanto:
$$\gamma_{0l} = \frac{\chi_{0l}}{R}, \quad \lambda_{ml}^2 = \left( \frac{\chi_{0l}}{R} \right)^2 + \left( \frac{m\pi}{L} \right)^2.$$
\end{itemize}

\textit{3. Superposición y determinación de coeficientes}
Por el principio de superposición para ecuaciones lineales (Alonso, Sec. 2.3.3), la solución general es una doble serie sobre los modos axiales y radiales:
$$T(\rho, z, t) = \sum_{m=1}^{\infty} \sum_{l=1}^{\infty} A_{ml} J_0\left( \frac{\chi_{0l}}{R} \rho \right) \sin\left( \frac{m\pi z}{L} \right) e^{-\alpha \lambda_{ml}^2 t}.$$
Aplicamos la condición inicial $T(\rho, z, 0) = T_0$:
$$T_0 = \sum_{m=1}^{\infty} \sum_{l=1}^{\infty} A_{ml} J_0\left( \frac{\chi_{0l}}{R} \rho \right) \sin\left( \frac{m\pi z}{L} \right).$$
Para despejar $A_{ml}$, explotamos la ortogonalidad de las bases de Fourier y de Bessel. Multiplicamos por $\sin(m'\pi z/L) J_0(\chi_{0l'} \rho/R) \rho$ e integramos en $z \in [0, L]$ y $\rho \in [0, R]$. Utilizando las relaciones de ortogonalidad (Alonso, Sec. 5.5.1 y 8.3.2):
$$\int_0^L \sin\left(\frac{m\pi z}{L}\right) \sin\left(\frac{m'\pi z}{L}\right) dz = \frac{L}{2} \delta_{mm'}, \quad \int_0^R \rho J_0\left(\frac{\chi_{0l} \rho}{R}\right) J_0\left(\frac{\chi_{0l'} \rho}{R}\right) d\rho = \frac{R^2}{2} \left[ J_1(\chi_{0l}) \right]^2 \delta_{ll'},$$
obtenemos:
$$A_{ml} = \frac{4}{L R^2 \left[ J_1(\chi_{0l}) \right]^2} \int_0^L \int_0^R T_0 \sin\left( \frac{m\pi z}{L} \right) J_0\left( \frac{\chi_{0l} \rho}{R} \right) \rho \, d\rho \, dz.$$
Las integrales se evalúan analíticamente:
$$\int_0^L \sin\left( \frac{m\pi z}{L} \right) dz = \frac{L}{m\pi} \left[ 1 - \cos(m\pi) \right] = \begin{cases} \frac{2L}{m\pi}, & m \text{ impar} \\ 0, & m \text{ par} \end{cases}$$
$$\int_0^R \rho J_0\left( \frac{\chi_{0l} \rho}{R} \right) d\rho = \frac{R^2}{\chi_{0l}} J_1(\chi_{0l}),$$
donde se utilizó la identidad $\int x J_0(x) dx = x J_1(x)$. Sustituyendo:
$$A_{ml} = \frac{4 T_0}{m\pi \chi_{0l} J_1(\chi_{0l})}, \quad m = 1, 3, 5, \dots; \ l = 1, 2, 3, \dots$$
La distribución de temperatura es entonces:
$$T(\rho, z, t) = \sum_{m=1,3,\dots}^{\infty} \sum_{l=1}^{\infty} \frac{4 T_0}{m\pi \chi_{0l} J_1(\chi_{0l})} J_0\left( \frac{\chi_{0l} \rho}{R} \right) \sin\left( \frac{m\pi z}{L} \right) e^{-\alpha \left[ \left( \frac{\chi_{0l}}{R} \right)^2 + \left( \frac{m\pi}{L} \right)^2 \right] t}.$$

\textbf{Parte II: Superficie lateral a $T(R, z) = f(z)$ (Condiciones no homogéneas)}

Cuando la condición de frontera en $\rho = R$ es no homogénea, el método directo de separación de variables no es aplicable porque las autofunciones radiales ya no satisfacen $J_0(\gamma R) = 0$. La estrategia estándar, fundamentada en la teoría de Sturm-Liouville (Alonso, Sec. 6.3 y 6.7), consiste en descomponer la solución en una parte estacionaria (que satisface la condición no homogénea) y una parte transiente (que satisface condiciones homogéneas):
$$T(\rho, z, t) = T_{ss}(\rho, z) + u(\rho, z, t).$$

\textit{1. Estado estacionario $T_{ss}(\rho, z)$}
En régimen estacionario, $\partial T/\partial t = 0$, por lo que $T_{ss}$ satisface la ecuación de Laplace $\nabla^2 T_{ss} = 0$ con condiciones:
$$T_{ss}(R, z) = f(z), \quad T_{ss}(\rho, 0) = 0, \quad T_{ss}(\rho, L) = 0.$$
Separando variables $T_{ss} = P(\rho)Z(z)$, la ecuación axial sigue siendo $Z'' + k_m^2 Z = 0$ con $Z(0)=Z(L)=0$, dando $Z_m(z) = \sin(m\pi z/L)$. La ecuación radial se convierte en:
$$\frac{1}{\rho} \frac{d}{d\rho} \left( \rho \frac{dP}{d\rho} \right) - k_m^2 P = 0,$$
que es la ecuación de Bessel modificada de orden cero (Alonso, Sec. 8.3.7). Sus soluciones son $I_0(k_m \rho)$ y $K_0(k_m \rho)$. Finitud en $\rho=0$ exige descartar $K_0$. Por superposición:
$$T_{ss}(\rho, z) = \sum_{m=1}^{\infty} C_m \frac{I_0(k_m \rho)}{I_0(k_m R)} \sin(k_m z).$$
Aplicando $T_{ss}(R, z) = f(z)$ y usando la ortogonalidad de senos:
$$C_m = \frac{2}{L I_0(k_m R)} \int_0^L f(z) \sin\left( \frac{m\pi z}{L} \right) dz.$$

\textit{2. Parte transiente $u(\rho, z, t)$}
La función $u$ satisface la ecuación de calor con condiciones de frontera homogéneas en toda la superficie ($u=0$ en $\rho=R, z=0, z=L$) y condición inicial modificada:
$$u(\rho, z, 0) = T_0 - T_{ss}(\rho, z).$$
La forma funcional de $u$ es idéntica a la obtenida en la Parte I, pero con coeficientes $B_{ml}$ diferentes:
$$u(\rho, z, t) = \sum_{m=1}^{\infty} \sum_{l=1}^{\infty} B_{ml} J_0\left( \frac{\chi_{0l} \rho}{R} \right) \sin(k_m z) e^{-\alpha \lambda_{ml}^2 t}.$$
Los coeficientes $B_{ml}$ se determinan proyectando la condición inicial sobre la base ortonormal:
$$B_{ml} = \frac{4}{L R^2 J_1^2(\chi_{0l})} \int_0^L \int_0^R \left[ T_0 - T_{ss}(\rho, z) \right] J_0\left( \frac{\chi_{0l} \rho}{R} \right) \sin(k_m z) \rho \, d\rho \, dz.$$
Sustituyendo $T_{ss}$ y aprovechando la ortogonalidad, el término $T_0$ genera el mismo coeficiente $A_{ml}$ de la Parte I, mientras que el término $-T_{ss}$ genera una corrección acoplada que involucra integrales mixtas de Bessel ordinarias y modificadas. El resultado final es:
$$T(\rho, z, t) = \sum_{m=1}^{\infty} C_m \frac{I_0(k_m \rho)}{I_0(k_m R)} \sin(k_m z) + \sum_{m=1}^{\infty} \sum_{l=1}^{\infty} B_{ml} J_0\left( \frac{\chi_{0l} \rho}{R} \right) \sin(k_m z) e^{-\alpha \lambda_{ml}^2 t}.$$

\textbf{Discusión física y convergencia}
La solución obtenida refleja la estructura de autovalores y autofunciones propia de problemas de Sturm-Liouville en dominios cilíndricos (Alonso, Sec. 6.3 y 8.3.2). El término exponencial $e^{-\alpha \lambda_{ml}^2 t}$ garantiza que los modos de alta frecuencia espacial (grandes $m$ y $l$) se amortiguan rápidamente, dominando el comportamiento a largo plazo el modo fundamental $(m=1, l=1)$. La convergencia de las series está asegurada por el decaimiento factorial de los coeficientes de Bessel y por la suavidad de las condiciones iniciales y de frontera. En el caso no homogéneo, la presencia de $I_0(k_m \rho)$ en el estado estacionario muestra cómo el calor penetra radialmente desde la superficie lateral hacia el eje, mientras que la parte transiente ajusta dinámicamente el perfil desde la condición inicial uniforme hasta el perfil de equilibrio impuesto por $f(z)$. Este formalismo es directamente aplicable a problemas de ingeniería térmica, diseño de reactores y análisis de procesos de templado en geometrías cilíndricas.


\section{Problema 9}
Una partícula de masa $m$ y sin espín se encuentra confinada en un cilindro circular de radio $R$ y longitud $L$. La ecuación de movimiento asociada a dicha partícula está dada por la ecuación de Schrödinger independiente del tiempo y sin potencial
$$-\frac{\hbar^2}{2m}\nabla^2\psi = E\psi.$$
Asumiendo que la función de onda $\psi$ se anula en las paredes, muestre que los autovalores están dados por
$$E_{njl} = \frac{\hbar^2}{2m}\left(\frac{\chi_{nj}^2}{R^2} + \frac{l^2\pi^2}{L^2}\right),$$
donde $\chi_{nj}$ son los ceros de la función de Bessel $J_n$ y $l=1,2,3,\dots$.

\subsection{Solución}
Para resolver este problema, partimos de la ecuación de Schrödinger independiente del tiempo para una partícula libre en el interior de una región de potencial nulo, tal como se introduce en la Sección 4.7 del texto \textit{Lecciones de física matemática} de Alonso Sepúlveda. La ecuación gobernante es:
$$-\frac{\hbar^2}{2m}\nabla^2\psi(\mathbf{r}) = E\psi(\mathbf{r}),$$
la cual puede reescribirse en la forma de la ecuación de Helmholtz homogénea:
$$\nabla^2\psi(\mathbf{r}) + k^2\psi(\mathbf{r}) = 0, \quad \text{con } k^2 = \frac{2mE}{\hbar^2}.$$
Dado que la geometría del confinamiento es cilíndrica, resulta natural expresar el laplaciano en coordenadas cilíndricas $(\rho, \phi, z)$, cuya forma explícita se detalla en el Anexo 1.1 del libro de Alonso:
$$\frac{1}{\rho}\frac{\partial}{\partial\rho}\left(\rho\frac{\partial\psi}{\partial\rho}\right) + \frac{1}{\rho^2}\frac{\partial^2\psi}{\partial\phi^2} + \frac{\partial^2\psi}{\partial z^2} + k^2\psi = 0.$$
Las condiciones de frontera del problema corresponden a un pozo de potencial infinito en las paredes del cilindro, lo que exige que la función de onda se anule en la superficie límite:
$$\psi(\rho=R, \phi, z) = 0, \quad \psi(\rho, \phi, z=0) = 0, \quad \psi(\rho, \phi, z=L) = 0.$$
Adicionalmente, la regularidad en el eje $\rho=0$ y la periodicidad azimutal $\psi(\rho, \phi+2\pi, z) = \psi(\rho, \phi, z)$ imponen restricciones topológicas y analíticas fundamentales.

\textbf{1. Separación de variables}
Aplicamos el método de separación de variables, ampliamente desarrollado en la Sección 3.3.3 de Alonso, proponiendo una solución factorizada:
$$\psi(\rho, \phi, z) = P(\rho)\,\Phi(\phi)\,Z(z).$$
Sustituyendo en la ecuación de Helmholtz y dividiendo por $P\Phi Z$, obtenemos:
$$\frac{1}{\rho P}\frac{d}{d\rho}\left(\rho\frac{dP}{d\rho}\right) + \frac{1}{\rho^2\Phi}\frac{d^2\Phi}{d\phi^2} + \frac{1}{Z}\frac{d^2Z}{dz^2} + k^2 = 0.$$
Multiplicando por $\rho^2$ para aislar la dependencia angular:
$$\frac{\rho}{P}\frac{d}{d\rho}\left(\rho\frac{dP}{d\rho}\right) + \frac{1}{\Phi}\frac{d^2\Phi}{d\phi^2} + \frac{\rho^2}{Z}\frac{d^2Z}{dz^2} + k^2\rho^2 = 0.$$
Dado que el segundo término depende exclusivamente de $\phi$, debe igualarse a una constante de separación. Por conveniencia matemática y física, la elegimos como $-n^2$, donde $n$ será un número entero:
$$\frac{1}{\Phi}\frac{d^2\Phi}{d\phi^2} = -n^2 \implies \frac{d^2\Phi}{d\phi^2} + n^2\Phi = 0.$$
Sustituyendo y reagrupando los términos restantes:
$$\frac{\rho}{P}\frac{d}{d\rho}\left(\rho\frac{dP}{d\rho}\right) - n^2 + \frac{\rho^2}{Z}\frac{d^2Z}{dz^2} + k^2\rho^2 = 0 \implies \frac{1}{Z}\frac{d^2Z}{dz^2} = -k_z^2,$$
donde $-k_z^2$ es la segunda constante de separación. Finalmente, la parte radial satisface:
$$\frac{1}{\rho}\frac{d}{d\rho}\left(\rho\frac{dP}{d\rho}\right) + \left(k^2 - k_z^2 - \frac{n^2}{\rho^2}\right)P = 0.$$
Definiendo $k_\rho^2 = k^2 - k_z^2$, obtenemos tres ecuaciones diferenciales ordinarias desacopladas.

\textbf{2. Solución de la parte axial y cuantización de $k_z$}
La ecuación para $Z(z)$ es:
$$\frac{d^2Z}{dz^2} + k_z^2 Z = 0.$$
Su solución general es $Z(z) = A\sin(k_z z) + B\cos(k_z z)$. Aplicando las condiciones de frontera de Dirichlet homogéneas en las bases del cilindro (Sección 2.3.3 de Alonso):
$$Z(0) = 0 \implies B = 0, \quad Z(L) = 0 \implies \sin(k_z L) = 0.$$
Esto exige que $k_z L = l\pi$ con $l = 1, 2, 3, \dots$ (el valor $l=0$ daría la solución trivial nula). Por tanto:
$$k_{z,l} = \frac{l\pi}{L}, \quad Z_l(z) = \sin\left(\frac{l\pi z}{L}\right).$$

\textbf{3. Solución de la parte azimutal y cuantización de $n$}
La ecuación angular $\Phi'' + n^2\Phi = 0$ tiene solución $\Phi(\phi) = C e^{in\phi} + D e^{-in\phi}$. La condición de unicidad y continuidad de la función de onda al dar una vuelta completa alrededor del eje $z$ exige $\Phi(\phi) = \Phi(\phi+2\pi)$, lo que impone que $n$ debe ser un número entero: $n = 0, \pm 1, \pm 2, \dots$. Esta cuantización topológica es fundamental en la mecánica cuántica de sistemas con simetría rotacional, como se discute en la Sección 3.3.3 y 8.4.5 de Alonso.

\textbf{4. Solución de la parte radial y aparición de las funciones de Bessel}
Sustituyendo $k_\rho^2$ en la ecuación radial obtenemos:
$$\frac{d^2P}{d\rho^2} + \frac{1}{\rho}\frac{dP}{d\rho} + \left(k_\rho^2 - \frac{n^2}{\rho^2}\right)P = 0.$$
Realizando el cambio de variable adimensional $x = k_\rho\rho$, esta ecuación se transforma exactamente en la ecuación de Bessel de orden $n$, estudiada rigurosamente en la Sección 8.3 del texto de Alonso:
$$x^2\frac{d^2P}{dx^2} + x\frac{dP}{dx} + (x^2 - n^2)P = 0.$$
La solución general es una combinación lineal de las funciones de Bessel de primera y segunda especie (Neumann): $P(x) = C J_n(x) + D N_n(x)$. Dado que la partícula puede ocupar todo el volumen del cilindro, incluyendo el eje $\rho=0$, y la función de onda debe ser finita en todo punto del dominio físico, debemos descartar $N_n(x)$ porque diverge logarítmicamente o como potencia negativa en el origen (Propiedad 4, Sección 8.3.1 de Alonso). Por tanto, $D=0$ y:
$$P(\rho) = J_n(k_\rho\rho).$$
Aplicando ahora la condición de frontera en la pared lateral $\rho = R$:
$$P(R) = 0 \implies J_n(k_\rho R) = 0.$$
Esto significa que $k_\rho R$ debe coincidir con uno de los ceros positivos de la función de Bessel de orden $n$. Denotando estos ceros como $\chi_{nj}$, donde $j=1,2,3,\dots$ numera el $j$-ésimo cero de $J_n$ (Sección 8.3.1, Propiedad 1 y Tabla de ceros), obtenemos la cuantización radial:
$$k_{\rho,nj} = \frac{\chi_{nj}}{R}.$$
Cada par $(n,j)$ define un modo radial permitido, con autofunción $J_n(\chi_{nj}\rho/R)$. La ortogonalidad y completez de esta base discreta de Bessel, fundamental para la expansión de estados cuánticos arbitrarios, se garantiza por la teoría de Sturm-Liouville desarrollada en la Sección 8.3.2 de Alonso.

\textbf{5. Espectro de energías (autovalores)}
Recordando la relación entre el número de onda total y la energía: $k^2 = \frac{2mE}{\hbar^2}$, y la descomposición espectral $k^2 = k_\rho^2 + k_z^2$, sustituimos los valores cuantizados obtenidos:
$$\frac{2mE_{njl}}{\hbar^2} = \left(\frac{\chi_{nj}}{R}\right)^2 + \left(\frac{l\pi}{L}\right)^2.$$
Despejando $E_{njl}$, llegamos exactamente a la expresión solicitada:
$$E_{njl} = \frac{\hbar^2}{2m}\left(\frac{\chi_{nj}^2}{R^2} + \frac{l^2\pi^2}{L^2}\right).$$
Los tres números cuánticos $(n,j,l)$ etiquetan de manera única cada estado estacionario: $n$ determina el momento angular azimutal, $j$ cuenta los nodos radiales (ceros de $J_n$), y $l$ cuenta los nodos axiales. La degeneración energética, típica de sistemas con simetría continua, ocurre cuando distintos tripletes $(n,j,l)$ producen el mismo valor de $\chi_{nj}^2/R^2 + l^2\pi^2/L^2$, fenómeno analizado en la Sección 6.5 de Alonso.

\textbf{Discusión física y conclusión}
La solución obtenida revela cómo la geometría cilíndrica impone una discretización de los modos espaciales permitidos, traduciéndose directamente en un espectro de energías cuantizadas. El término $\chi_{nj}^2/R^2$ proviene del confinamiento radial y refleja la estructura oscilatoria de las funciones de Bessel, mientras que $l^2\pi^2/L^2$ corresponde al confinamiento axial y reproduce el espectro de una partícula en una caja unidimensional. La condición de anulación en las paredes (potencial infinito) garantiza que el operador Hamiltoniano sea hermítico bajo las condiciones de frontera de Dirichlet homogéneas (Sección 6.2.4 de Alonso), asegurando autovalores reales y autofunciones ortonormales. Este formalismo constituye la base para el estudio de nanoestructuras cilíndricas, guías de onda cuánticas, y estados electrónicos en nanotubos o puntos cuánticos de geometría alargada, donde la separación de variables en coordenadas cilíndricas y la teoría espectral de Bessel son herramientas analíticas insustituibles.


\bibliographystyle{plainurl}
\bibliography{bibliografia} 
\end{document}
